\chapter{Review of Related Literature}
\label{ch:literature}

Activated sludge surrogate modeling draws on several bodies of research that address different parts of the same engineering problem. Mechanistic models define the treatment processes, component transformations, and material balances that a surrogate must approximate. Data-driven studies show how treatment quantities can be estimated from measurements or simulations. Physical enforcement methods constrain predictions to satisfy specified balances and concentration bounds. Interpretable modeling provides ways to examine the fitted relationships, while optimization studies determine whether those predictions can support useful operating decisions. This chapter reviews these contributions as a connected body of work and identifies the questions that remain when they are brought together for complete-component activated sludge analysis.

\section{Mechanistic foundations of activated sludge modeling}

\subsection{Development of component and nutrient models}

Activated sludge modeling developed from the need to relate operating conditions to organic matter removal, oxygen demand, biomass growth, and sludge production. Early formulations established the basic connection between biological transformation and material accounting. \citet{Marais1976} formulated steady-state activated sludge analysis around substrate utilization and biomass, while \citet{Dold1981} developed a more general model that linked several biological processes within a common treatment calculation. Together, these contributions established the basis for calculating treatment responses from material balances and biological rates rather than treating effluent quality as an isolated empirical outcome.

This framework was expanded as activated sludge models became more component-resolved. \citet{Henze1987} presented the single-sludge framework associated with Activated Sludge Model No.\ 1 (ASM1), which distinguishes substrate, biomass, and nitrogen components so that carbon oxidation, nitrification, and denitrification can be calculated together. Later models addressed processes that this formulation did not represent explicitly. \citet{Henze1999} extended phosphorus-removal modeling to include the denitrifying activity of phosphorus-accumulating organisms, while \citet{Gujer1999} introduced Activated Sludge Model No.\ 3 (ASM3), which represents organic substrate storage and uses endogenous respiration to describe biomass decay.

These developments addressed different biological modeling needs. The phosphorus-removal extension connects nitrate use with phosphate behavior, while the storage-based formulation separates substrate uptake from its later use by biomass. The collected formulations in \citet{Henze2006} provide a common reference for these component definitions and their associated reaction terms. For surrogate modeling, their importance lies in defining what must be predicted. Effluent COD or TN can summarize treatment quality, but the underlying component concentrations describe the material forms that produce those totals. A surrogate intended to preserve process accounting therefore requires more than agreement with a small set of aggregate indicators.

The comparison of seven activated sludge models by \citet{Hauduc2013} further shows why the mechanistic reference must be specified before statistical models are compared. Their review examined how the models represent hydrolysis, fermentation, microbial growth and decay, and biological storage. The differences were not merely mathematical; they reflected different assumptions about treatment processes and available process knowledge. Model selection therefore depends on the treatment system and the question being studied. For a surrogate comparison, this means that all candidate predictors must approximate the same selected process model. Otherwise, differences in prediction may reflect different biological assumptions rather than differences among statistical learners.

\subsection{Biological dependence and calibration uncertainty}

The mechanistic literature also shows why operating conditions and influent composition cannot be interpreted independently. Nutrient-removal processes depend on the material entering the system and on how operating choices alter the conditions under which microorganisms compete and react. \citet{Wentzel1992} examined interactions among nitrification, denitrification, and excess phosphorus removal. \citet{Guerrero2011} studied competition for organic substrate in a pilot system designed for simultaneous nitrogen and phosphorus removal and found that the supplied carbon source affected the ability to maintain phosphorus removal when nitrate entered the initially anaerobic zone. Their results illustrate a specific biological dependence among substrate composition, recycle conditions, and nutrient removal.

For statistical modeling, these findings support retaining both operating conditions and component-resolved influent information. A model based only on total organic loading can overlook differences among carbon sources that perform different biological roles. Likewise, a model that treats aeration or recycle independently of loading can miss the conditions under which those operating changes influence nutrient removal. This provides a process-based reason to examine interaction terms in a treatment surrogate. It does not, however, imply that a fitted interaction coefficient identifies a microbial mechanism by itself. The statistical relationship and the biological mechanism remain distinct forms of evidence.

Calibration research addresses a related problem: even a mechanistically detailed model requires evidence that its parameters and internal quantities represent the treatment system of interest. \citet{Petersen2003} reviewed experimental designs for estimating activated sludge parameters, while \citet{Vanrolleghem1999} examined the use of respirometry, which measures oxygen consumption, to estimate model parameters and component combinations. These methods show that the information supplied by an experiment must correspond to the parameter or internal quantity being estimated.

The full-scale calibration study of \citet{Machado2014} makes this issue more explicit. They calibrated a phosphorus-removal model using data from a treatment plant in Manresa, Spain, assessed parameter confidence intervals, and examined how influent and operating variables affected the model fit. They also varied these inputs while retaining default kinetic and stoichiometric parameters to evaluate output sensitivity. Their analysis therefore considered uncertainty in both process inputs and biological parameters. For surrogate modeling, this distinction is important because simulated training targets inherit the assumptions of the mechanistic model that generates them. Agreement with those targets establishes how well the surrogate approximates that reference. Agreement with an operating treatment plant requires the additional evidence supplied by calibration and measurement.

The mechanistic literature therefore establishes the component definitions, reaction relationships, and calibration practices needed to give a surrogate prediction physical meaning. The unresolved question for this dissertation is not how to replace that mechanistic foundation, but how to approximate its defined treatment responses efficiently while retaining their material accounting. The surrogate is evaluated against an adopted mechanistic reference; calibration of that reference to a particular operating plant remains a separate task.

\subsection{Clarification and connected plant balances}

Biological reactors alone do not determine the final treatment response. The material leaving the reactor is subsequently separated, recycled, wasted, or discharged, so clarification and solids storage become part of the plant-level material account. \citet{Takacs1991} developed a clarification and thickening model based on solids fluxes and balances across layers of a one-dimensional settler. The model calculates solids profiles and overflow and underflow concentrations under steady and changing conditions. \citet{HartelPopel1992} likewise modeled secondary clarification with sludge thickening. These studies provide mechanistic descriptions of how suspended material is retained, concentrated, recycled, and discharged.

The difficulty of connecting biological and clarification models was examined directly by \citet{JeppssonDiehl1996}. They studied how biological components propagate through a secondary clarifier and identified problems that arise because the biological reactor and settler models represent particulate material differently. Their evaluation of component-transport algorithms focused particularly on the composition of sludge returned to the reactor. This result is important for surrogate modeling because it shows that predicting aggregate solids alone is not sufficient when the recycled stream must retain a biologically meaningful component composition.

\citet{Alex2008} addressed the broader need for comparable whole-plant assessment through Benchmark Simulation Model No.\ 1. The benchmark connects five biological compartments, a secondary settler, internal recycle, return sludge, and wasting, and specifies influent conditions, control settings, and assessment criteria. Its contribution is not only a plant configuration but a common setting in which alternative operating strategies can be compared. It makes explicit that biological conversion, solids separation, recirculation, and wasting form one connected treatment response.

Together, these studies define the relationships that a connected surrogate must preserve. A clarifier prediction affects both the discharged effluent and the recycled biomass returned to the biological tanks. Stored solids contribute to the inventory used in solids-retention calculations. Shared streams carry material from one unit to another and must therefore agree across unit boundaries. A single-reactor surrogate captures only part of this system. Extending physical enforcement to the plant level requires consistent component flows at these connections together with an explicit account of retained solids.

\section{Data-driven prediction of wastewater treatment responses}

\subsection{Soft sensors and reported treatment quantities}

Data-driven wastewater treatment research developed partly because operating plants collect large sets of measurements that are difficult to interpret jointly and because some important treatment quantities are costly, delayed, or difficult to measure directly. \citet{Newhart2019} reviewed methods for fault detection, variable prediction, and advanced process control and identified missing observations, outliers, correlated variables, and changing process behavior as recurring features of plant data. \citet{Sundui2021} reviewed machine learning applications across biological wastewater treatment. Together, these reviews establish a range of useful prediction tasks and emphasize that the choice of method depends on both the available information and the intended application.

Soft sensors are one example. A soft sensor estimates a quantity that is difficult or slow to measure from other observations that are easier to obtain. \citet{Ching2022} developed a soft sensor for five-day biochemical oxygen demand using extreme gradient boosting (XGBoost). They evaluated influent and effluent estimates at two treatment plants and used supporting measurements selected from variables such as COD, solids, nutrients, and simpler sensor readings. The model covered a wide range of oxygen-demand values and could accommodate missing supporting observations. Its contribution is therefore practical and specific: it provides a faster estimate of an otherwise slow water-quality measurement.

A similar distinction appears in effluent nutrient prediction. \citet{Wang2025} compared a mechanistic activated sludge model with six machine learning models for effluent TN prediction using one year of high-resolution municipal plant data. The mechanistic model reproduced the direction of TN variation, while the learned models achieved greater predictive accuracy in the assessed dataset. Feature-attribution analysis was also used to identify influential inputs. This study shows that statistical models can improve a defined effluent-prediction task while still allowing the fitted response to be inspected.

Other data-driven studies address operational quantities beyond effluent quality. \citet{Ekinci2023} predicted waste-sludge production using 208 daily observations from an advanced biological treatment plant in Turkey. Their inputs included flow, polymer dosage, and the amounts of solids and pollutants removed. They compared learning methods and reported improved prediction after feature selection. This adds a solids-management target to the quality indicators considered elsewhere. At the same time, its use of measured removal quantities illustrates how the available inputs depend on the prediction task. A model used before a proposed operating condition has been evaluated would require inputs available at that earlier decision stage.

The common achievement of these studies is reliable estimation of selected treatment quantities from available plant information. Such predictions can support monitoring, diagnosis, and particular operating analyses. The requirement in this dissertation is broader. When the goal is to check mass conservation, non-negativity, and the distribution of material among biological forms, the complete modeled component state becomes the relevant prediction target. Accuracy for TN, oxygen demand, or sludge production supplies useful evidence for those particular quantities, but it does not establish the accuracy or physical consistency of every modeled component.

\subsection{Component estimation and simulation-based surrogates}

Component-estimation studies move closer to the information used by mechanistic activated sludge models. \citet{Wang2024} used a backpropagation neural network to estimate five COD-related component quantities from measured influent COD. They first characterized those components in an operating plant and then evaluated the neural estimates against the measurements. Their reported prediction accuracy exceeded 80\% for the five assessed components. The study addresses a practical difficulty in mechanistic modeling: repeated experimental fractionation of influent COD components is demanding, yet those fractions are needed as model inputs.

That contribution differs from predicting the effluent component state produced by treatment. Influent fractionation estimates how incoming organic material is distributed among model categories. An effluent surrogate instead estimates how a defined influent changes under specified biological and operating conditions. The two tasks can complement one another, but they have different targets and therefore require different evidence. The component estimates of \citet{Wang2024} improve characterization of the material entering the model; they do not by themselves establish the complete treatment mapping investigated in this dissertation.

Simulation-based surrogates address that mapping more directly. \citet{Fang2011} combined an extended activated sludge model with a support vector surrogate and a genetic optimization search for a full-scale nutrient-removal system. Mechanistic simulations supplied the input--output relationship learned by the surrogate. Their optimization indicated that anoxic volume and internal recycle could be reduced while retaining the specified discharge quality. The study is an early example of a workflow in which mechanistic simulation generates the training data, statistical learning accelerates response evaluation, and optimization uses the learned approximation to select operating conditions.

\citet{He2023} used a similar simulation-to-surrogate logic for greenhouse-gas emissions from papermaking wastewater treatment. They fitted a Kriging surrogate to steady-state process simulations and reported a coefficient of determination above 0.96. The surrogate then supported analysis and optimization of controllable treatment variables. This work addresses the computational cost of repeatedly evaluating an environmental objective. Its prediction target, however, is the emissions measure rather than the complete component state required for the material-balance assessment of this dissertation.

A closer precedent is provided by \citet{Selerio2026}, who addressed complete single-reactor component prediction using invariant-constrained second-order regression. The study used 10,000 simulated steady-state cases with 20 effluent components and compared the structured regression with nine statistical predictors. Its assessment included changes in training size, errors in reported water-quality quantities, and physical violation diagnostics. This establishes a direct component-level activated sludge comparison and provides the foundation for the physical-enforcement and interpretation questions developed later in this chapter.

Table~\ref{tab:rrl_prediction_evidence} summarizes this progression from monitoring selected indicators to approximating complete process responses. The distinction is not that one task is inherently better than another. Rather, each prediction target establishes a different kind of evidence. The table therefore identifies what additional information is needed when the objective shifts from estimating a reported quantity to evaluating a complete activated sludge component state.

\begin{table}[htbp]
\centering\footnotesize
\caption{Treatment prediction studies and the scope of their evidence.}
\label{tab:rrl_prediction_evidence}
\begin{tabular}{>{\raggedright\arraybackslash}p{0.24\linewidth}>{\raggedright\arraybackslash}p{0.33\linewidth}>{\raggedright\arraybackslash}p{0.31\linewidth}}
\toprule Study and task & Contribution established & Additional question for this dissertation\\\midrule
\citet{Ching2022}\newline Oxygen-demand soft sensing & Estimation across two plants with support for missing sensor observations. & Can a predictor estimate the complete effluent state from influent and proposed operating inputs?\\
\citet{Wang2025}\newline Effluent TN prediction & Comparison of six learned models with a mechanistic model and inspection of influential inputs. & Are the individual nitrogen and other component concentrations accurate and physically consistent?\\
\citet{Wang2024}\newline Influent component estimation & Neural estimation of five COD-related quantities used in process modeling. & How accurately is the resulting effluent component state predicted after treatment?\\
\citet{Fang2011}\newline Simulation-based operating search & Mechanistic simulation combined with a support vector surrogate and optimization. & Can connected component predictions be constrained and checked at the selected controls?\\
\citet{Selerio2026}\newline Complete reactor prediction & Structured prediction of 20 components with physical checks and training-size assessment. & How do common projection and plant connections affect prediction and operating assessment?\\
\bottomrule
\end{tabular}
\end{table}

\subsection{Model comparison, data availability, and extrapolation}

The wastewater applications reviewed above use statistical learners with different assumptions about how nonlinear input--output relationships should be represented. Random forests, developed by \citet{Breiman2001}, average predictions from randomized decision trees. \citet{ChenGuestrin2016} developed XGBoost, while \citet{Ke2017} introduced the light gradient boosting machine (LightGBM). \citet{SmolaScholkopf2004} described support vector regression, which combines regularized fitting with kernel-based similarity. These methodological studies provide the statistical basis for several learners now used in wastewater treatment. Their differences justify comparing multiple model families rather than assuming that one fitting structure is best for every treatment response.

Inspectable linear and nonlinear models provide another part of that comparison. \citet{Wold2001} described partial least squares regression, which represents related variation in inputs and outputs through latent components. These latent statistical components are weighted combinations of variables and should not be confused with physical wastewater constituents. \citet{Woo2009} applied a kernel extension of partial least squares to estimate COD, TN, and cyanide concentrations in industrial cokes wastewater treatment. The nonlinear extension outperformed the assessed linear and other nonlinear partial least squares formulations in that application. This supports the value of nonlinear structure for the studied response while leaving ordinary partial least squares as a distinct baseline for other comparison tasks.

Broader tabular-learning research reinforces the value of comparing model families. \citet{ShwartzZiv2022} examined tabular machine learning, while \citet{Borisov2024} surveyed neural approaches for tabular data. Activated sludge surrogate inputs have a similar tabular form, with each case containing operating settings and influent concentrations. However, these methodological comparisons cannot determine the preferred learner for a specific treatment problem. That ranking must come from a common treatment target, a common data design, and a common evaluation procedure. The general literature therefore guides comparator selection, while the activated sludge benchmark supplies the process-specific evidence.

The amount of available training data adds another dimension to the comparison. \citet{VieringLoog2023} reviewed learning curves and showed that their shapes depend on the learner, the dataset, and the prediction task. A learning curve relates model performance to the number of training observations and therefore provides evidence about data efficiency. \citet{Selerio2026} applied this type of assessment to activated sludge reactor surrogates. The result characterizes performance within the stated input range, but a broader comparison still needs to determine how complete-state error and physical adjustment change together when the number of simulated training cases becomes small.

Such comparisons also depend on strict separation between model development and assessment information. \citet{Kaufman2012} defined information leakage and described how information unavailable at the intended prediction time can inadvertently enter model construction. In wastewater applications, this distinction is particularly important because downstream measurements may be available for monitoring but not for predicting a proposed operating condition before the treatment response has occurred. Scaling, preprocessing, and model tuning must likewise exclude observations reserved for final assessment. These requirements determine what an error score means; they do not imply that downstream measurements are inappropriate in studies whose stated purpose is monitoring.

The remaining comparative question therefore combines three issues: learner choice, training support, and operating domain. Existing indicator and component studies demonstrate that useful prediction is possible within their assessed conditions. What is still needed here is a common complete-state comparison in which several learners use the same inputs, outputs, data partitions, and error definitions. That assessment must also distinguish interpolation within the sampled domain from prediction beyond it. Only then can component error and physical validity be evaluated together as influent and operating conditions change.

\section{Methods for imposing physical constraints on predictions}

\subsection{Physical guidance during fitting}

Machine learning models can reproduce observed responses while violating the equations or bounds that define a physically acceptable process state. A major branch of physics-informed learning addresses this problem by incorporating governing relations directly into model fitting. \citet{Karniadakis2021} reviewed approaches that combine physical knowledge with machine learning. One influential example is the residual-based formulation of \citet{Raissi2019}, in which discrepancies in governing differential equations are added to the training objective. Equation disagreement is therefore penalized alongside prediction error, allowing process knowledge to guide learning even when response observations are limited.

Wastewater applications provide direct evidence for this idea. \citet{Li2024} examined physics-informed neural models for water-treatment unit dynamics, including an activated sludge reactor. They compared physically informed models with conventional learned models under limited data and changing loading conditions. The physically informed models showed stronger predictive capability and generalization in the assessed systems. Their results demonstrate that governing equations can provide useful learning information when data alone are insufficient.

The distinction between physical guidance and physical enforcement is important, however. A finite residual penalty balances equation agreement against the other terms in the fitting objective. It can reduce violations without requiring every new prediction to satisfy the relevant equality exactly. The contribution of \citet{Li2024} therefore concerns improved learning through physical information. If the engineering requirement is that every returned component state satisfy mass conservation and concentration bounds, an additional output construction, projection, or numerical check is still required. Physical guidance during fitting and exact sample-specific enforcement address complementary questions.

\subsection{Conservation through model construction}

A second approach is to build conservation directly into the way model outputs are constructed. \citet{Beucler2021} introduced neural architectures that encode analytic constraints and compared them with constraints imposed through a loss function. In their climate-convection application, the architectural formulation enforced the selected conservation laws to machine precision without degrading overall predictive performance. Their results show that, when the output representation is designed appropriately, known equalities can be satisfied by construction rather than encouraged through a penalty.

Related work has used stoichiometric structure to couple predicted species. \citet{SturmWexler2020} developed a conservation framework in which learned transfers are converted into species changes through balanced process relationships. \citet{SturmWexler2022} later placed the stoichiometric mapping in a fixed neural layer for a photochemical surrogate. This construction avoided the need to reconstruct all reaction transfers as separate learning targets while retaining atom conservation in the predicted species changes. The common idea is that conservation is carried by the mapping between process changes and the final species response.

That idea transfers naturally to activated sludge modeling because stoichiometry also defines allowable combinations of component changes. However, conservation alone does not guarantee a physically acceptable concentration state. A balanced change can still consume more of a component than is available and thereby produce a negative concentration. The distinction follows from the difference between constraining a reaction-induced change and constraining the final state. The physical claim must therefore be matched to the quantity that the model construction actually restricts.

\citet{Harder2022} illustrated this broader requirement in an aerosol microphysics surrogate by examining both physical penalties and constrained output completion. Their assessment treated mass conservation and positive mass as separate but simultaneous concerns. The lesson for activated sludge prediction is direct: the complete returned state must be assessed for both balance satisfaction and concentration bounds after every stage of prediction has been applied.

\subsection{Projection after statistical prediction}

Projection provides another route to physical enforcement. Instead of changing the fitted learner, a prediction is first produced and then adjusted to satisfy specified constraints. \citet{SturmSilva2024} developed a species-weighted least-squares correction for atom conservation in an atmospheric chemistry surrogate. Their weights accounted for the magnitude and uncertainty of individual species changes. In the reported assessment, an unweighted correction could damage prediction of a small radical, whereas weighting retained its accuracy and slightly improved overall agreement. The enforced conservation relation was unchanged; what changed was the way the required adjustment was distributed among species.

The same issue arises in activated sludge prediction because component concentrations can differ by several orders of magnitude. A correction that is numerically small relative to an abundant substrate may be large relative to a residual nutrient concentration. The results of \citet{SturmSilva2024} therefore motivate assessing projection size and component error together rather than treating physical correction as neutral with respect to prediction quality. Their method addresses an equality constraint, however, so an additional construction is needed when non-negativity must also be guaranteed.

\citet{KircherVotsmeier2025} addressed this joint requirement through a positivity-preserving projection followed by linear interpolation backtracking. They evaluated the method on atmospheric chemistry, heterogeneous catalysis, and synthetic reaction networks. The corrected predictions satisfied atom balances and remained positive in the tested systems without reducing overall prediction accuracy. This establishes a general procedure for imposing both conservation and positivity across different underlying surrogates. The remaining question for activated sludge analysis is how such a procedure behaves with wastewater component scales, reactor boundaries, limited training data, and inputs that move beyond the sampled domain.

A direct activated sludge example is provided by \citet{Selerio2026}. The structured regression in that study used physical penalties during fitting, then checked each raw prediction and projected it onto the mass conservation equalities. If non-negativity still failed, a linear program---an optimization problem with a linear objective and linear constraints---was used to complete the correction. The published assessment reported no final mass conservation or non-negativity violations. Because the raw regression itself did not satisfy both requirements, the final enforcement procedure was the stage responsible for the physical guarantee.

That result also illustrates why physical validity and predictive accuracy must be evaluated separately. The structured surrogate was not the most accurate method under every reported score, even though its final states satisfied the declared constraints. Physical enforcement therefore establishes that the returned state belongs to the specified feasible set; it does not establish that the state is the closest possible approximation to the mechanistic reference. The next comparative question is consequently a paired one: when the same raw predictions are held fixed, how does projection change error, adjustment size, and computational cost across several activated sludge learners?

Figure~\ref{fig:rrl_enforcement_locations} places the reviewed methods along the prediction pipeline. Fitting penalties guide learning toward equation agreement, constrained output constructions impose selected relations through the model architecture, and final-state projection corrects the prediction after it has been produced. These approaches can all introduce physical information, but they make different guarantees because they act at different stages.

\begin{figure}[htbp]
\centering
\begingroup\linespread{1}\selectfont
\begin{tikzpicture}[>=Latex,node distance=5mm,
heading/.style={draw=black!75,rounded corners,fill=cyan!6,text width=4.1cm,minimum height=8mm,align=center,font=\footnotesize,inner sep=4pt},
block/.style={heading,fill=white,minimum height=17mm},
limit/.style={block,fill=black!5},
flow/.style={->,thick,draw=black!75}]
\node[heading] (lh) at (-4.55,0) {Construction};
\node[heading] (mh) at (0,0) {Training};
\node[heading] (rh) at (4.55,0) {Returned state};
\node[block,below=of lh] (l1) {Model process states\\and physical relations};
\node[block,below=of mh] (m1) {Fit observations with\\physical residual penalties};
\node[block,below=of rh] (r1) {Apply a defined\\projection to a\\completed prediction};
\node[block,below=of l1] (l2) {Enforce encoded relations\\within the\\output mapping};
\node[block,below=of m1] (m2) {Penalties guide the fit\\alongside the data objective};
\node[block,below=of r1] (r2) {Impose declared relations\\on each returned state};
\node[limit,below=of l2] (l3) {Guarantee depends on\\the stated model};
\node[limit,below=of m2] (m3) {Small residuals do not\\certify every new output};
\node[limit,below=of r2] (r3) {Feasibility alone leaves\\process accuracy\\unresolved};
\node[heading,text width=13.3cm,minimum height=12mm,below=7mm of m3] (common)
{Compare mass conservation and non-negativity on the same returned state};
\draw[flow] (lh)--(l1); \draw[flow] (l1)--(l2); \draw[flow] (l2)--(l3);
\draw[flow] (mh)--(m1); \draw[flow] (m1)--(m2); \draw[flow] (m2)--(m3);
\draw[flow] (rh)--(r1); \draw[flow] (r1)--(r2); \draw[flow] (r2)--(r3);
\draw[flow] (l3.south)--(common.north west);
\draw[flow] (m3.south)--(common.north);
\draw[flow] (r3.south)--(common.north east);
\end{tikzpicture}

\endgroup
\caption{Stages at which the reviewed methods introduce physical information. Fitting penalties guide equation agreement. Output construction and projection impose the relations encoded in their formulations. Simultaneous mass conservation and non-negativity require an assessment of the final component state.}
\label{fig:rrl_enforcement_locations}
\end{figure}

\subsection{Constraint scope and recent extensions}

Recent work has expanded constrained learning beyond deterministic point predictions. \citet{Hansen2024} used integral conservation relations to adjust probabilistic predictions while accounting for both predicted means and uncertainty. Their framework enforced conservation exactly when the imposed balance was assigned zero uncertainty and was evaluated on a family of transport equations. \citet{Utkarsh2025} trained a probabilistic predictor together with a differentiable projection and evaluated the approach on differential-equation and time-series problems. Because the projection is differentiable, its effect can participate directly in gradient-based learning. These studies extend physical enforcement from correcting one predicted state to representing uncertainty within the constrained prediction process.

Other work has examined the expressive capacity of hard-constrained models. \citet{Min2024} developed neural output layers for input-dependent linear and convex constraints and established approximation results under the stated mathematical conditions. A convex feasible set contains the straight line joining any two feasible states, which permits useful mathematical guarantees for constrained outputs. \citet{Valente2025} examined projection onto physical manifolds, or sets of outputs satisfying selected physical laws, using spring-mass and reactive-plasma examples. Their work extends correction beyond simple linear relations to nonlinear physical structures.

These contributions broaden the available tools for physically constrained machine learning, but they do not remove the need to define the wastewater-specific constraint set carefully. The literature already provides methods for enforcing conservation and, in some cases, conservation together with positivity. The gap addressed in this dissertation is therefore not the absence of physical enforcement methods. It is the need to determine how common enforcement affects accuracy, adjustment size, and computational cost for a complete activated sludge component response and then to extend the material-accounting requirements from a single reactor to a connected plant. The imposed constraints remain intentionally narrower than the complete biological and settling equations.

\section{Interpretable modeling of treatment responses}

\subsection{Inspection of inputs and fitted relationships}

Interpretability concerns what can be learned about a model's response after it has been fitted. \citet{Rudin2019} argued for using directly inspectable models in consequential applications when suitable interpretable alternatives are available. In wastewater treatment, this becomes relevant when an engineer needs to understand how influent composition, aeration, recycle, or another operating variable contributes to a predicted treatment outcome. The argument supports examining the model that informs the decision. It does not imply that every visible statistical coefficient has a direct biological interpretation.

Several treatment studies use post-fitting inspection to identify influential variables. \citet{Wang2025} combined effluent TN prediction with feature-attribution analysis, while \citet{Ekinci2023} used feature selection and attribution in sludge prediction. These approaches help identify which observed variables contribute to a particular statistical output and can guide process investigation or reduce unnecessary inputs. Their explanations, however, apply to the fitted TN or sludge model. They do not describe how every component in a mechanistic activated sludge state responds.

Explicit algebraic models provide a complementary form of interpretability. \citet{Nazlabadi2021} reviewed response-surface methods in biological wastewater treatment, where polynomial terms describe operating effects and interactions over a fitted domain. \citet{Brunton2016} developed sparse equation discovery from candidate functions for nonlinear dynamical systems. The latter provides a methodological precedent for inspecting named terms, while response-surface applications provide a direct wastewater-treatment context. Together, they show why visible mathematical structure can be useful when the objective is to understand how inputs contribute to a predicted response.

For activated sludge analysis, the distinction among terms matters. An aeration term represents a fitted dependence on aeration, an influent substrate term represents a fitted dependence on loading, and their product allows the aeration response to change with the substrate condition. Such terms can make the statistical response easier to examine within the fitted domain. They do not, by themselves, establish causal biological parameters. Demonstrating a microbial mechanism would require a different identification problem and different evidence, particularly when the model inputs are correlated.

\subsection{Structured regression and the final projected state}

\citet{Selerio2026} addressed the need for an inspectable complete-component activated sludge surrogate by separating operating effects, influent effects, curvature, interactions, and learned dependence among effluent components. The explicit second-order terms allow the fitted input relationships to be examined in physical coordinates. Component coupling then contributes to the complete raw state, and a fixed composition map converts that state into reported water-quality quantities. The result is a treatment surrogate whose statistical structure can be traced from inputs to component predictions and then to aggregate treatment indicators.

The same study also shows why interpretability cannot stop at the fitted coefficients when the deployed prediction includes physical correction. Its displayed coefficient relationships describe the raw regression before final constraint enforcement. Mass conservation projection and the subsequent linear-program stage can change the concentrations ultimately returned to the engineer. If a concentration bound becomes active, the local sensitivity of the final prediction can also differ from the slope implied by the unconstrained equation. The study therefore distinguishes the displayed raw coefficient structure from derivatives of the final projected response.

This distinction becomes especially important when component coupling is present. A visible coefficient can affect one intermediate quantity, which then influences several components through the coupling relation before projection imposes the physical constraints. Different fitted coefficient combinations can also produce similar predictions. \citet{Selerio2026} therefore identifies coefficient stability across resampled training data as an additional diagnostic rather than assuming that one fitted coefficient set uniquely represents the treatment process.

The broader literature on differentiable hybrid modeling supports this complete-response perspective. \citet{Shen2023} reviewed approaches that connect learned relationships with process equations and showed, through geoscience examples, how combined models can retain physical knowledge while remaining open to interpretation. They also identify verification of the physical significance of outputs as an ongoing need. For the present dissertation, this supports following the prediction through all of its stages rather than interpreting only the first fitted equation.

The remaining interpretive question is therefore an extension of an established contribution. The raw structured regression already exposes its terms. What remains is to connect those terms to component coupling, reporting, and the constraints that determine the final effluent response, and to examine how stable the inferred relationships remain as training cases change. This preserves the distinction between statistical association and biological mechanism while making the complete deployed calculation available for engineering interpretation.

\section{Surrogate modeling for operating decisions}

\subsection{Mechanistic and learned routes to operating optimization}

Optimization of activated sludge operation predates the recent use of machine learning surrogates. Mechanistic studies established how treatment constraints, influent conditions, and operating degrees of freedom can be brought into one decision problem. \citet{Araujo2013} optimized a steady-state process model and used sensitivity analysis to identify variables for economic control. Their analysis linked active effluent constraints and remaining operating degrees of freedom to the resulting control structure. \citet{Nazif2023} calibrated a treatment model to a plant in Tehran, grouped influents according to quantitative and qualitative characteristics, and optimized operating settings for each group. These studies provide examples of direct process-based optimization in which operating decisions are evaluated through the mechanistic treatment model itself.

Learned surrogates provide an alternative when repeated mechanistic calculations are expensive. \citet{Fang2011} trained a support vector surrogate on mechanistic simulations and then used the learned model in an optimization of a nutrient-removal system. \citet{KusiakWei2013} fitted neural models for airflow, effluent carbonaceous biochemical oxygen demand, and TSS from industrial data and used a multiobjective particle-swarm search to examine dissolved-oxygen settings. These studies demonstrate that a learned treatment response can be used directly in the search for operating conditions. The engineering meaning of the result, however, depends on the particular outputs that the surrogate predicts and on how those outputs represent the objectives and constraints of the decision.

Dynamic control studies provide related evidence at a different time scale. \citet{Bernardelli2020} developed a neuro-fuzzy predictive controller for aeration and field-tested it at a large municipal plant, reporting improved effluent quality and energy savings in that application. \citet{Croll2023} reviewed reinforcement learning for wastewater control and discussed the methods and challenges involved in changing operation over time. These contributions show that learned models can participate directly in treatment control. The present dissertation addresses a different question: fixed-input steady-state operating assessment. Dynamic control results therefore provide important context without replacing the need for a separate evaluation of the steady-state decision problem.

The operating literature establishes that both mechanistic and learned routes can be used to select treatment settings. The unresolved question for a new surrogate is whether its selected settings remain favorable when both routes are judged on the same engineering basis. Such a comparison requires common treatment objectives, operating bounds, and reference calculations. It also requires separating the one-time cost of developing a surrogate from the computational cost of repeatedly using it during an operating search.

\subsection{Integrated resource recovery and connected responses}

Whole-system wastewater studies broaden the decision beyond the performance of a single biological reactor. \citet{PadronPaez2020} examined sustainable treatment design through multiple objectives, while \citet{Aboagye2021} combined wastewater-network design, optimization, and sustainability assessment. Their work shows how treatment quality and resource priorities can be considered together when several process choices are available. This provides the broader systems context for a connected plant objective in which treatment gains must be considered alongside the resources required to obtain them.

\citet{Durkin2024} extended this idea through surrogate-based process synthesis for resource recovery from brewery wastewater. Their framework combined high-fidelity simulations, regression surrogates, and classification models for process feasibility. The application considered biogas recovery, recoverable nutrient quality, aeration demand, and uncertainty in wastewater composition. It therefore demonstrates that learned approximations can support integrated decisions involving several resource pathways while accounting explicitly for whether a simulated configuration is feasible.

The decision problem in this dissertation is related but different. \citet{Durkin2024} considered a network-design problem with alternative process configurations. The present connected operating study retains a specified activated sludge topology and predicts the mixer, biological reactor stages, clarifier outlets, and stored solids jointly. Its additional requirement is that those predicted quantities remain mutually consistent at the fixed physical connections of the plant. A classifier indicating simulator feasibility and an explicitly imposed component balance therefore address different aspects of model reliability.

The single-reactor work of \citet{Selerio2026} supplies a precedent for physical enforcement at one biological unit, while the benchmark of \citet{Alex2008} provides a mechanistic precedent for comparing connected plant responses and controls. Bringing these lines of research together leads to a specific remaining question: can a statistical plant response preserve the declared unit connections and solids inventories while still supporting useful operating decisions? Answering that question requires evidence beyond single-reactor feasibility and beyond the accuracy of one aggregate operating objective.

\subsection{Approximation error at selected operating conditions}

A surrogate can have good average predictive accuracy and still be inaccurate where the optimizer ultimately chooses to operate. Surrogate-optimization research addresses this distinction directly. \citet{BhosekarIerapetritou2018} reviewed surrogate modeling for feasibility analysis and optimization, while \citet{McBride2019} reviewed surrogate construction from experiments and simulations for process-system decisions. Both place prediction within a larger decision workflow: the surrogate must approximate the quantities and the operating region that matter to the optimization, not merely reproduce an average response over the original training design.

\citet{Alexandrov1998} developed a trust-region framework that manages approximation models during optimization by using reference information to assess local model quality and control the region in which the approximation is trusted. The contribution is methodological, but its implication for wastewater treatment is clear. A model used to select aeration, recycle, or wasting settings should be checked near the controls that the optimizer actually selects rather than judged only by an average error calculated elsewhere in the operating domain.

The distinction between predictive quality and optimization quality was examined explicitly by \citet{Pedrozo2025}. They compared five surrogate approaches for a carbon-dioxide-stream pooling problem and assessed predictive accuracy, computational effort, and optimization behavior separately. Their comparison of one-shot optimization with a trust-region filter strategy showed that models with favorable fitting performance could still differ in convergence and reliability during search. Although the application concerns carbon capture rather than wastewater treatment, the methodological lesson is directly relevant: a good surrogate fit and a good engineering decision are separate outcomes.

Physical feasibility does not eliminate this distinction. Mass conservation and non-negativity remove certain impossible component states, but many different physically admissible states can still have different nitrogen distributions, biomass concentrations, and treatment performance. An optimizer can therefore select a prediction that satisfies the declared physical constraints yet remains inaccurate relative to the mechanistic reference. Physical-enforcement studies establish how to exclude impossible states; approximation-management studies establish why the response used by the optimizer still needs to be checked.

Together, these findings motivate independent mechanistic evaluation of the aeration, recycle, wasting, and other controls selected by the surrogate. The remaining decision gap concerns the complete body of evidence for a connected activated sludge surrogate. Its unit balances and stored solids must be mutually consistent. Its selected controls must satisfy the original treatment requirements when those controls are returned to the reference model. Its objective must be evaluated on the same mechanistic basis as the direct comparator. Search time and verification effort must also be reported on a clearly defined basis. Only then can computational benefit and decision quality be assessed without assuming that physical constraint satisfaction establishes either one.

\section{Synthesis of established contributions and remaining gaps}

\subsection{What the reviewed studies have addressed}

The literature reviewed in this chapter shows a clear progression from process representation to engineering decision support. Mechanistic activated sludge models define components, biological reactions, separation, recirculation, and solids inventories. Calibration studies connect those model quantities to observations from real treatment systems. Data-driven soft sensors and indicator models improve estimation of selected treatment responses. Simulation-based surrogates extend learned prediction to repeated process calculations and optimization. Physics-informed and constrained-learning methods introduce governing relations into fitting or impose them on model outputs. Interpretable regression and surrogate-optimization research then address how fitted responses can be examined and how they should be evaluated when used to make operating decisions.

Several requirements relevant to this dissertation have therefore already been addressed in important forms. \citet{KircherVotsmeier2025} established joint atom balance and positivity for specified reactive-chemistry surrogates. \citet{Selerio2026} established final mass conservation and non-negativity for a structured activated sludge reactor surrogate and showed why interpretation must distinguish the raw and projected responses. \citet{Durkin2024} demonstrated surrogate-supported optimization for integrated wastewater resource recovery. These studies are not gaps that the dissertation claims to fill. They are the starting points from which its more specific questions arise.

The remaining gaps appear when these established contributions are brought into one complete-component activated sludge assessment. A comparison of learners must use the same component target. Physical enforcement must be evaluated not only for feasibility but also for the errors and adjustments it introduces. Interpretation must follow the prediction through coupling and projection rather than stopping at the visible regression terms. Plant-level decisions must preserve unit connections and then be checked on the original mechanistic basis. Table~\ref{tab:rrl_evidence_gap_matrix} summarizes this relationship between the established literature and the specific questions retained for the dissertation.

\begin{table}[htbp]
\centering\footnotesize
\caption{Established studies and the remaining questions addressed by the dissertation.}
\label{tab:rrl_evidence_gap_matrix}
\begin{tabular}{>{\raggedright\arraybackslash}p{0.22\linewidth}>{\raggedright\arraybackslash}p{0.34\linewidth}>{\raggedright\arraybackslash}p{0.32\linewidth}}
\toprule Research need & Studies addressing the need & Remaining question\\\midrule
Complete treatment prediction & \citet{Wang2025} assesses TN prediction. \citet{Wang2024} estimates influent components. \citet{Selerio2026} predicts complete reactor components. & How do several learners compare on one complete-state target under training scarcity and extrapolation?\\
Physical validity of predictions & \citet{Beucler2021} and \citet{SturmWexler2022} impose conservation. \citet{KircherVotsmeier2025} and \citet{Selerio2026} address conservation and concentration bounds together. & How does common enforcement affect component accuracy, adjustment size, and cost across learners?\\
Interpretation of treatment responses & \citet{Wang2025} inspects input contributions. \citet{Selerio2026} exposes regression terms and distinguishes raw from final interpretation. & How do coupling, reporting, and active constraints affect sensitivities and stability of the explanation?\\
Connected operating decisions & \citet{Fang2011} and \citet{Durkin2024} use surrogates in treatment decisions. \citet{Alex2008} defines a connected plant benchmark. & Can joint component enforcement support settings that remain credible under independent plant evaluation?\\
\bottomrule
\end{tabular}
\end{table}

\subsection{The four linked questions of the dissertation}

The first gap concerns comparative prediction of the complete reactor response. Existing indicator studies show that machine learning can predict useful treatment quantities, while the structured-regression study of \citet{Selerio2026} extends prediction to all modeled reactor components. What remains is a common comparison across learning families in which the target, partitions, and assessment scales are held fixed. That comparison must also show how performance changes as the number of training cases decreases and as influent or operating conditions move beyond the sampled domain. The purpose is not to identify one universally superior learner, but to determine how learner choice depends on the available data and the intended input range.

The second gap concerns what physical enforcement changes after a raw prediction has already been produced. Joint mass conservation and concentration bounds have been established by existing methods. The remaining question is how a common enforcement procedure affects the same raw predictions across several activated sludge learners. Component error, physical residuals, negative concentrations, projection displacement, and computational cost describe different consequences of that correction. Evaluating them together makes it possible to determine the value of enforcement without assuming that physical compliance must also improve predictive accuracy.

The third gap concerns interpretation of the complete returned treatment state. Structured regression already provides visible operating, loading, curvature, interaction, and coupling terms. However, the response ultimately used in engineering analysis can differ from the raw equation because coupling, reporting, and projection occur afterward. The remaining question is therefore how sensitivities pass through those stages, particularly when concentration bounds become active, and whether the resulting fitted relationships remain stable when the training cases change. This extends interpretation from the visible regression to the complete prediction actually returned to the engineer.

The fourth gap concerns connected prediction and operating assessment. Mechanistic plant models establish the shared streams, clarification behavior, recycle structure, and stored solids that connect treatment units. Surrogate studies establish that learned models can support operating and resource-recovery decisions. What remains here is to combine those requirements in one connected component response, impose the declared plant relations jointly, and then independently calculate treatment outcomes at the controls selected by the surrogate. Comparison with direct mechanistic optimization places both routes on the same treatment and resource basis and separates faster search from better decision quality.

Figure~\ref{fig:rrl_intersecting_gaps} shows why these four questions cannot be treated independently. The component concentrations being predicted are also the quantities adjusted by physical enforcement and examined during interpretation. Those same concentrations determine treatment indicators and enter the operating objective and constraints. A change introduced by projection can therefore affect component error, local sensitivity, and the settings selected during optimization. The integrated assessment must follow this shared treatment response through each stage rather than allowing evidence from one stage to stand in for another.

\begin{figure}[htbp]
\centering
\begingroup\linespread{1}\selectfont
\begin{tikzpicture}[>=Latex,node distance=7mm,
block/.style={draw=black!75,rounded corners,fill=cyan!6,align=center,font=\footnotesize,inner sep=4pt},
evidence/.style={block,text width=4.1cm,minimum height=20mm,fill=white},
flow/.style={->,thick,draw=black!75}]
\node[block,text width=12.9cm,minimum height=12mm] (scope) at (0,0)
{Defined process response and decision-time information};
\node[evidence] (stat) at (-4.55,-2.25)
{Component comparison\\Prediction accuracy\\Training size\\Extrapolation};
\node[evidence] (phys) at (0,-2.25)
{Physical assessment\\Mass conservation\\Non-negativity\\Accuracy and cost};
\node[evidence] (interp) at (4.55,-2.25)
{Final prediction analysis\\Coupling and projection\\Local sensitivity\\Coefficient stability};
\node[block,text width=10.1cm,minimum height=18mm,below=8mm of phys] (decision)
{Connected operating assessment\\Shared streams and stored solids\\Mechanistic evaluation at selected controls};
\node[block,text width=12.9cm,minimum height=14mm,fill=black!5,below=7mm of decision] (gap)
{Common research focus\\Assessment of the same learned and projected treatment responses};
\draw[flow] (scope.south west)--(stat.north);
\draw[flow] (scope.south)--(phys.north);
\draw[flow] (scope.south east)--(interp.north);
\draw[flow] (stat.south)--(decision.north west);
\draw[flow] (phys.south)--(decision.north);
\draw[flow] (interp.south)--(decision.north east);
\draw[flow] (decision)--(gap);
\end{tikzpicture}

\endgroup
\caption{Links among the remaining comparative, physical, interpretive, and operating questions. Each question concerns the same defined activated sludge response. Established methods address parts of the problem, while their combined use requires common component and decision assessment.}
\label{fig:rrl_intersecting_gaps}
\end{figure}

The literature therefore supports a progression from a defined reactor response to a physically checked and independently evaluated plant decision. Chapter~\ref{ch:foundations} develops the common component and mechanistic basis for that progression. Chapter~\ref{ch:projection} compares reactor predictors and examines the effect of physical enforcement on the same raw responses. Chapter~\ref{ch:interpretable} develops the structured regression and follows interpretation through component coupling and projection. Chapter~\ref{ch:plant} extends the framework to a connected activated sludge plant and independently evaluates the controls selected through surrogate-assisted and direct mechanistic optimization. Chapters~\ref{ch:projection}--\ref{ch:plant} each present the rationale, theory, methodology, results, and discussion for their corresponding assessment, while Chapter~\ref{ch:conclusion} brings their findings together and identifies directions for future research.